\begin{statementbox}

	\textbf{\textcolor{CStatement}{条件}}
	\begin{description}[style=nextline, leftmargin=2.8em, labelsep=0.5em, itemsep=0.35em]
		\item[\textbf{\textcolor{CStatement}{条件 1（变号因子）：}}] $k(x)$ 在区间 $D$ 上变号（即 $k(x)$ 既取正值也取负值）。
		\item[\textbf{\textcolor{CStatement}{条件 2（定义区间）：}}] $D$ 是一个区间，且函数 $k(x), h(x)$ 在 $D$ 上有定义。
	\end{description}

	\textbf{\textcolor{CStatement}{结论}}
	\begin{description}[style=nextline, leftmargin=2.8em, labelsep=0.5em, itemsep=0.45em]
		\item[\textbf{\textcolor{CStatement}{结论一（等价分类法）：}}]
			\textbf{\textcolor{CStatement}{直观描述：}}
			按 $k(x)$ 的正负零将 $D$ 分为三段，分别化为 $a$ 与 $h/k$ 的不等式，再取交集得到参数范围。
			其中 $D_1=\{x\in D\mid k(x)>0\}$，$D_2=\{x\in D\mid k(x)<0\}$，$D_0=\{x\in D\mid k(x)=0\}$。
			\[
				\begin{cases}
					a\ge \dfrac{h(x)}{k(x)}, & \forall x\in D_1,\\[6pt]
					a\le \dfrac{h(x)}{k(x)}, & \forall x\in D_2,\\[6pt]
					0\ge h(x), & \forall x\in D_0.
			\end{cases}
			\]

		\item[\textbf{\textcolor{CStatement}{常见变形(不等号反向):}}]
			\textbf{\textcolor{CStatement}{直观描述:}}
			若原不等式为 $a\cdot k(x)\le h(x)$,则只需将正负区间的不等号方向反转.
			\[
				\begin{cases}
					a\le \dfrac{h(x)}{k(x)}, & \forall x\in D_1,\\[6pt]
					a\ge \dfrac{h(x)}{k(x)}, & \forall x\in D_2,\\[6pt]
					0\le h(x), & \forall x\in D_0.
			\end{cases}
			\]
	\end{description}

\end{statementbox}
